\documentclass[11pt]{article}
\usepackage[a4paper,margin=1in]{geometry}
\usepackage{amsmath,amsthm,amssymb}
\usepackage{mathtools}

\newtheorem{theorem}{Theorem}[section]
\newtheorem{proposition}[theorem]{Proposition}
\newtheorem{lemma}[theorem]{Lemma}
\newtheorem{corollary}[theorem]{Corollary}
\theoremstyle{definition}
\newtheorem{definition}[theorem]{Definition}
\newtheorem{remark}[theorem]{Remark}

\newcommand{\R}{\mathbb{R}}
\newcommand{\Fr}{\mathcal{A}}
\newcommand{\Om}{\Omega}
\newcommand{\eps}{\varepsilon}
\newcommand{\rstar}{\rho_{*}}
\newcommand{\zbar}{\bar z_{\rho}}
\newcommand{\er}{e_{r,\zbar}}
\newcommand{\PiE}{\Pi(E_{\rho},\zbar)}

\title{Space-Time $\Pi$ Identity and the Integrated\\
       Quadratic Boundary Oscillation\\
       (Route~$\delta$, Plan~2, Building Block \textsc{SpaceTimeIdentity})}
\author{}
\date{}

\begin{document}
\maketitle

\begin{abstract}
This auxiliary block records the exact space-time radial-moment identity
\[
   \int_{\rstar}^{\rho_{\delta}}\PiE\,d\mu(\rho)
   =(n-1)\!\int_{\widetilde{\Om}}\!\sigma(x)\,dx,
\]
and isolates the quantitative centroid-route estimate
\[
   \int_{\rstar}^{\rho_\delta}\PiE\,d\mu(\rho)
   \le C_\Pi(n,R,\rstar)\,\delta_T(\Om).
\tag{$\Pi$-int}
\]
Conditional on \textup{($\Pi$-int)}, we prove the integrated
quadratic boundary oscillation
\[
   \int_{\rstar}^{\rho_{\delta}}\!\!\int_{\partial^{*}E_{\rho}}|\nu_{E_{\rho}}-\er|^{2}\,d\mathcal H^{n-1}\,d\mu(\rho)
   \le C(n,R,\rstar)\,\delta_{T}(\Om),
   \quad\textup{(B$^{2}$-int)}
\]
and, under an additional annular/perimeter cap for the centroid slices, the
integrated radial trace
\[
   \int_{\rstar}^{\rho_{\delta}}\!\!\Bigl(\int_{\partial^{*}E_{\rho}}|H_{\zbar,\rho}-1|\,d\mathcal H^{n-1}\Bigr)^{2}d\mu(\rho)
   \le C(n,R,\rstar)\,\delta_{T}(\Om),
   \quad\textup{(T$_{1}^{2}$-int)}
\]
The proof of (T$_{1}^{2}$-int) uses a
\emph{slicing-then-squaring} decomposition $T_{1}=T_{1}^{\rm rad,slic}+T_{1}^{\rm tan}$,
in which the radial piece is bounded \emph{linearly} (not at the FMP
rate $\sqrt{\delta}$) via the good-ray equality of slicewise volumes
and the bad-ray multiplicity bound of \cite{LevelSetIdentity}; the
squared volume Fraenkel $|E_{\rho}\Delta B_{\rho}|^{2}$ is controlled
at rate $\delta$ via the new pointwise slice-rearrangement inequality
$|E_{\rho}\Delta B_{\rho}|^{2}\le C\,\PiE$. The repaired downstream
metric route in \cite{WeightedMetricTrace} no longer depends on
\textup{($\Pi$-int)}, but this block remains a useful diagnostic for the
centroid-based space--time approach.
The invalid shortcut from the rearranged Talenti profile to
\textup{($\Pi$-int)} has been removed; proving \textup{($\Pi$-int)}
would be an additional structural theorem, not a prerequisite for the
current metric closure.
\end{abstract}

\section{Notation and statement of the targets}\label{sec:targets}

\subsection{Standing hypotheses}

Fix $n\ge 2$, $R\ge 2$, $\rstar\in(0,1)$. Let $\Om\subset B_R$ open
with $|\Om|=\omega_n$, $\delta:=\delta_T(\Om)\le\delta_0(n,R,\rstar)$.
Let $u=u_\Om$ be the variational torsion function,
$t:[\rstar,1]\to(0,\|u\|_\infty)$ the volume-radius parametrisation,
$E_\rho:=\{u>t(\rho)\}$, $t_B(\rho):=(1-\rho^2)/(2n)$,
$\psi(\rho):=(-t_B'(\rho))-(-t'(\rho))=\rho/n+t'(\rho)\ge 0$,
$\rho_\delta:=1-\kappa\sqrt\delta$,
$\mu(d\rho):=(-t'(\rho))d\rho$.
By coarea applied to $u\in W^{1,2}_0$, $E_\rho$ has finite
perimeter for a.e.\ $\rho$.

For each $\rho$, $\zbar=\bar z_\rho$ is the centroid of $E_\rho$
(\cite{CentroidBound} Theorem F). Since $E_\rho\subset B_R$, also
$\zbar\in B_R$.  No Cicalese--Leonardi containment is assumed at this
stage; any lower bound on $|x-\zbar|$ must be justified separately on
good levels or by truncating the kernel.

The $L^2$ divergence identity in \cite{LevelSetIdentity}:
\begin{equation}\label{eq:53prime}
   \int_{\partial^*E_\rho}|\nu_{E_\rho}-\er|^2 d\mathcal H^{n-1}
   = 2(P(E_\rho)-P(B_\rho))+2\PiE,
\end{equation}
where
\begin{equation}\label{eq:Pi}
   \PiE := (n-1)\int_{B_\rho(\zbar)\setminus E_\rho}|x-\zbar|^{-1}dx
         - (n-1)\int_{E_\rho\setminus B_\rho(\zbar)}|x-\zbar|^{-1}dx \ge 0.
\end{equation}

\subsection{The three target theorems}

\begin{theorem}[Space-time $\Pi$ identity, (G1)]\label{thm:G1}
$\int_{\rstar}^{\rho_\delta}\PiE d\mu(\rho)=(n-1)\int_{\widetilde\Om}\sigma(x)dx$
for a signed Fubini density $\sigma$ satisfying the kernel-weighted bound
\[
   |\sigma(x)|\le \int_{J_x}|x-\bar z_\rho|^{-1}\,d\mu(\rho),
\]
where
$J_x:=\{\rho:x\in B_\rho(\zbar)\Delta E_\rho\}$.
\end{theorem}

\begin{theorem}[Integrated (B$^2$-int) from \textup{($\Pi$-int)}]\label{thm:G3}
Assume \textup{($\Pi$-int)}:
\[
   \int_{\rstar}^{\rho_\delta}\PiE\,d\mu(\rho)
   \le C_\Pi(n,R,\rstar)\,\delta_T(\Om).
\]
Then
\[
   \int_{\rstar}^{\rho_\delta}\int_{\partial^*E_\rho}|\nu_{E_\rho}-\er|^2 d\mathcal H^{n-1} d\mu(\rho)
   \le C_{B^2}(n,R,\rstar)\,\delta_T(\Om).
\]
\end{theorem}

\begin{theorem}[Auxiliary integrated radial trace, (T$_1^2$-int)]\label{thm:T1sq}
Assume \textup{($\Pi$-int)} and the auxiliary annular/perimeter cap
\[
   \rstar/2\le |x-\bar z_\rho|\le 2R \text{ on the relevant boundary pieces},
   \qquad
   P(E_\rho)-P(B_\rho)+\PiE\le M_A(n,R,\rstar)
\]
for a.e.\ $\rho\in[\rstar,\rho_\delta]$. Then
\[
   \int_{\rstar}^{\rho_\delta}\Bigl(\int_{\partial^*E_\rho}|H_{\zbar,\rho}-1|d\mathcal H^{n-1}\Bigr)^2 d\mu(\rho)
   \le C_{\rm radial}(n,R,\rstar)\,\delta_T(\Om).
\]
\end{theorem}

\section{Space-time identity (G1)}\label{sec:G1}

By \eqref{eq:Pi}, $\PiE = (n-1)\int_{\R^n}(\mathbf 1_{B_\rho(\zbar)}-\mathbf 1_{E_\rho})|x-\zbar|^{-1}dx$.
By joint Borel measurability and Fubini,
\begin{equation}\label{eq:fubini}
   \int_{\rstar}^{\rho_\delta}\PiE d\mu(\rho)
   = (n-1)\int_{\R^n}\sigma(x)dx,
\end{equation}
where $\sigma(x):=\int_{\rstar}^{\rho_\delta}(\mathbf 1_{B_\rho(\zbar)}-\mathbf 1_{E_\rho})/|x-\zbar|\,d\mu(\rho)$.
The integrand vanishes outside
$\widetilde\Om:=B_{R+1}$ (which contains $\Om\subset B_R$), because $\zbar\in B_R$ and
$\rho\le1$. The kernel is locally integrable in dimension $n\ge2$; if desired,
one first proves the identity with $\min\{|x-\zbar|^{-1},M\}$ and then lets
$M\to\infty$ by monotone convergence applied to the positive and negative
parts. Taking absolute values inside the $\rho$-integral gives the stated
kernel-weighted crossing bound. \hfill\qed

\section{The remaining quantitative input}\label{sec:G2}

The identity of Theorem~\ref{thm:G1} rewrites
$\int\PiE\,d\mu$ as an integral over the space--time mismatch
\[
   J_x:=\{\rho\in[\rstar,\rho_\delta]:
          x\in B_\rho(\zbar)\Delta E_\rho\}.
\]
The rearranged Talenti/Nicola--Tilli profile gap controls the
one-dimensional distributional profiles \(t(\rho)\) and \(t_B(\rho)\).
It does \emph{not}, by itself, control the physical crossing set
\(J_x\), because \(J_x\) also depends on the moving spatial radius
\(|x-\bar z_\rho|\). Therefore the implication
\[
   \mathcal L^1(J_x)\lesssim\delta_T(\Om)
\]
is not asserted here. The additional theorem that would finish this
centroid subroute is the following integrated moment estimate.

\begin{theorem}[Integrated \(\Pi\)-control; remaining centroid-route input]
\label{thm:PiInput}
There is \(C_\Pi(n,R,\rstar)>0\) such that, for every bounded open
\(\Om\subset B_R\) with \(|\Om|=\omega_n\) in the small-deficit regime,
\begin{equation}\label{eq:PiInput}
   \int_{\rstar}^{\rho_\delta}\PiE\,d\mu(\rho)
   \le C_\Pi(n,R,\rstar)\,\delta_T(\Om).
   \tag{\(\Pi\)-int}
\end{equation}
\end{theorem}

\begin{remark}[Possible proof route]
A viable proof of Theorem~\ref{thm:PiInput} would have to combine the
Fubini identity of Theorem~\ref{thm:G1} with a genuine space--time
coarea or transport estimate for the physical crossings of
\(|x-\bar z_\rho|=\rho\). The \(H^1\)-in-\(\rho\) centroid estimate of
\cite{CentroidBound} is the natural input for controlling the moving
centre term, but an extra argument is still needed to convert the
profile gap into the physical crossing-length bound.
\end{remark}

\section{Assembly of (B$^2$-int) from \textup{($\Pi$-int)}}\label{sec:G3}

Assume \eqref{eq:PiInput}. By \eqref{eq:53prime} integrated,
\[
   \int|\nu-\er|^2 dH d\mu = 2\int(P-P_\rho)d\mu+2\int\PiE d\mu.
\]
The first piece is $\le C\delta$ via the deficit identity in
\cite{LevelSetIdentity}.
The second piece is exactly \eqref{eq:PiInput}. Both pieces are
\(O(\delta)\), giving Theorem~\ref{thm:G3}. \hfill\qed

\section{The integrated radial trace (T$_1^2$-int) --- slicing-then-squaring}\label{sec:T1}

Write $T_1=T_1(E_\rho,\zbar):=\int_{\partial^*E_\rho}\lvert\,\lvert x-\zbar\rvert/\rho-1\,\rvert\,d\mathcal H^{n-1}$
and decompose
\[
   T_1 = T_1^{\rm rad,slic}+T_1^{\rm tan},
\]
with
\begin{align*}
   T_1^{\rm rad,slic}&:=\int_{\partial^*E_\rho}||x-\zbar|/\rho-1|\cdot|\er\cdot\nu_{E_\rho}|d\mathcal H^{n-1},\\
   T_1^{\rm tan}&:=\int_{\partial^*E_\rho}||x-\zbar|/\rho-1|\cdot(1-|\er\cdot\nu_{E_\rho}|)d\mathcal H^{n-1}.
\end{align*}

By the slicing identity \cite[Thm 18.11]{Maggi2012} applied with
$g(r):=|r/\rho-1|$:
\begin{equation}\label{eq:S1}
   T_1^{\rm rad,slic} = \int_{S^{n-1}}\sum_{j=1}^{N(\theta)}|r_{\theta,j}/\rho-1|\,r_{\theta,j}^{n-1}d\theta.
\end{equation}

\subsection{Pointwise bound for $T_1^{\rm rad,slic}$}

\begin{lemma}\label{lem:T1rad}
$T_1^{\rm rad,slic}\le C_4(n,R,\rstar)[|E_\rho\Delta B_\rho(\zbar)|+(P(E_\rho)-P(B_\rho))]$.
\end{lemma}

\begin{proof}
\emph{Good rays.} On a good ray ($N=1$), the slicewise 1D
symmetric-difference \emph{equals}:
\begin{equation}\label{eq:eq-slicewise}
   \int_0^\infty|\mathbf 1_{E_{\rho,\theta}}-\mathbf 1_{(0,\rho)}|r^{n-1}dr
   = \frac{1}{n}|r_{\theta,1}^n-\rho^n|.
\end{equation}
By MVT, $|r_{\theta,1}/\rho-1|\cdot r_{\theta,1}^{n-1}\le\beta_n|r_{\theta,1}^n-\rho^n|$
with $\beta_n=\beta_n(R,\rstar)$. Integrating and using spherical Fubini:
\[
   \int_{\{N=1\}}|r_{\theta,1}/\rho-1|r_{\theta,1}^{n-1}d\theta \le n\beta_n|E_\rho\Delta B_\rho(\zbar)|.
\]

\emph{Bad rays.} $|r_{\theta,j}/\rho-1|r_{\theta,j}^{n-1}\le C(R,\rstar)$,
summed gives $\le (2R)^{n-1}(1+k(\theta))$. We claim the bad-ray
multiplicity bound
\[
   \int_{S^{n-1}} k(\theta)\,d\theta
   \;\le\; (\rstar/2)^{-(n-1)}\,(P(E_\rho)-P(B_\rho)),
\]
which we prove as follows (note: this estimate is conditional on the
auxiliary annular cap of Theorem~\ref{thm:T1sq}, which provides
$|x-\zbar|\ge\rstar/2$ on the relevant boundary pieces). For each
direction $\theta\in S^{n-1}$, let $k(\theta)\ge 0$ count the crossings
of the radial ray with $\partial^*E_\rho$ beyond the spherical baseline.
By the radial-projection inequality and the area formula,
$\int_{S^{n-1}}\sum_j r_{\theta,j}^{n-1}\,d\theta = P_{\rm rad}(E_\rho)
\le P(E_\rho)$, where $r_{\theta,j}$ are the crossing radii. Each
crossing contributes $\ge(\rstar/2)^{n-1}$ to this integral, since
$|x-\zbar|\ge\rstar/2$ on the annular cap. Summing the contributions
beyond the spherical baseline gives
$\int_{S^{n-1}}k(\theta)\,d\theta \le (\rstar/2)^{-(n-1)}\,
(P(E_\rho)-P(B_\rho))$, as claimed.

Combining the two estimates gives the claim.
\end{proof}

\subsection{Pointwise bound for $T_1^{\rm tan}$}

\begin{lemma}\label{lem:T1tan}
$T_1^{\rm tan}\le (P(E_\rho)-P(B_\rho))+\PiE$.
\end{lemma}

\begin{proof}
Under the auxiliary annular cap, $||x-\zbar|/\rho-1|\le C(R,\rstar)$, so
$T_1^{\rm tan}\le C(R,\rstar)\int(1-|\er\cdot\nu|)dH\le C(R,\rstar)\int(1-\er\cdot\nu)dH=C(R,\rstar)((P-P_\rho)+\PiE)$
by the divergence identity in \cite{LevelSetIdentity}.
\end{proof}

\subsection{Squaring and integrating}

By Lemmas \ref{lem:T1rad}, \ref{lem:T1tan} and $(a+b+c)^2\le 3(a^2+b^2+c^2)$:
\[
   T_1^2 \le 3C_5^2[|E_\rho\Delta B_\rho|^2+(P-P_\rho)^2+\PiE^2].
\]
Integrating against $d\mu$ and bounding each piece by $C\delta$:
\begin{itemize}
\item $\int|E_\rho\Delta B_\rho|^2 d\mu\le C_8\int\PiE d\mu\le C\delta$
via the slice-rearrangement (§\ref{sec:slice}) + (G3).
\item $\int(P-P_\rho)^2 d\mu\le M_A\int(P-P_\rho)d\mu\le C\delta$
under the auxiliary cap (uses $(P-P_\rho)\ge 0$ to bound
$(P-P_\rho)^2\le M_A(P-P_\rho)$ on the cap).
\item $\int\PiE^2 d\mu\le M_A\int\PiE d\mu\le C\delta$ under the same cap.
\end{itemize}
Hence $\int T_1^2 d\mu\le C_{T_1^2}\delta$.

Combining with (G3) gives (T$_1^2$-int) as follows. Since
$H_{\zbar,\rho} = (|x-\zbar|/\rho)\,(\er\cdot\nu)$, we decompose
\[
   H_{\zbar,\rho} - 1 \;=\; \bigl(|x-\zbar|/\rho - 1\bigr)\,(\er\cdot\nu)
                       \;+\; \bigl(\er\cdot\nu - 1\bigr),
\]
where on the annular cap $|\er\cdot\nu|\le 1$ and
$\bigl||x-\zbar|/\rho-1\bigr|\le C(R,\rstar)$, so both factors are
bounded. Taking absolute values and integrating over $\partial^*E_\rho$
yields
\[
   \int_{\partial^*E_\rho}|H_{\zbar,\rho}-1|\,d\mathcal H^{n-1}
   \;\le\; T_1
   \;+\; \int_{\partial^*E_\rho}|\er\cdot\nu-1|\,d\mathcal H^{n-1}.
\]
For the second term, using $|\er\cdot\nu-1|\le|\er-\nu|$ and applying
Cauchy--Schwarz with respect to $d\mathcal H^{n-1}\,d\mu$ together with
the perimeter bound $P(E_\rho)\le M_A$ on the bounded reduction,
\[
   \Bigl(\int|\er-\nu|\,d\mathcal H^{n-1}\,d\mu\Bigr)^2
   \;\le\; M_A\cdot\int|\er-\nu|^2\,d\mathcal H^{n-1}\,d\mu.
\]
The $L^2$-bound on $|\er-\nu|^2$ provided by (G3) (Theorem~\ref{thm:G3})
then gives an $O(\delta)$ control. Squaring the pointwise decomposition,
applying $(a+b)^2\le 2(a^2+b^2)$, integrating against $d\mu$, and
combining with $\int T_1^2\,d\mu\le C_{T_1^2}\delta$ yields
(T$_1^2$-int). \hfill\qed

\section{The slice-rearrangement inequality}\label{sec:slice}

\begin{lemma}\label{lem:slice-rearrange}
Under the annular cap $\rstar/2\le r\le 2R$,
$|E_\rho\Delta B_\rho(\zbar)|^2\le C_8(n,R,\rstar)\PiE$ pointwise in
the smallness regime.
\end{lemma}

\begin{proof}
\emph{Step 1: $\PiE\ge c M_1$.} Let $M_1:=\int_{E_\rho\Delta B_\rho}|\rho-|x-\zbar||dx$.
By rewriting \eqref{eq:Pi},
\[
   \PiE = (n-1)\int_{E_\rho\Delta B_\rho}\frac{|\rho-|x-\zbar||}{\rho|x-\zbar|}dx \ge\frac{n-1}{2R\rho}M_1.
\]

\emph{Step 2: $|E_\rho\Delta B_\rho|^2\le C\,M_1$.} By polar coordinates
and a slicewise Cauchy--Schwarz on $[\rstar/2,2R]$,
$\sigma(\theta)^2\le 4R\,m_1(\theta)$ with
$\sigma(\theta)=\mathcal L^1(E_{\rho,\theta}\Delta(0,\rho))$,
$m_1(\theta)=\int|r-\rho|\mathbf 1_{E_{\rho,\theta}\Delta(0,\rho)}dr$.
Cauchy--Schwarz on $S^{n-1}$ then gives
$(\int\sigma d\theta)^2\le 4R|S^{n-1}|\int m_1 d\theta$. Since
$\sigma(\theta)\,d\theta\cdot r^{n-1}$ (for $\rstar/2\le r\le 2R$)
gives the volume element up to the comparable factors $(\rstar/2)^{n-1}$
and $(2R)^{n-1}$, the slice-to-volume conversion preserves the bound up
to constants depending only on $n,\rstar,R$; this yields
$|E_\rho\Delta B_\rho|^2\le C_9\,M_1$.

Combining Steps 1--2 gives the claim.
\end{proof}

\section{Discussion: why $T_1$-route closes where $T_2$-route stalls}\label{sec:discussion}

\begin{remark}
The $T_2$-route would extract on good rays
$(r_{\theta,1}/\rho-1)^2 r_{\theta,1}^{n-1}\le C|r_{\theta,1}^{n-1}-\rho^{n-1}|$
(a square on LHS, unsigned mismatch on RHS), then need
$\int|r_{\theta,1}^{n-1}-\rho^{n-1}|d\theta\le C\sqrt{\delta_T}$
(FMP rate). Integrating against $d\mu$ leaks at $\sqrt\delta$.

The $T_1$-route uses linear good-ray equality \eqref{eq:eq-slicewise}
and the slice-rearrangement to get $|E_\rho\Delta B_\rho|^2\le C\Pi$,
which integrates to $C\delta$ via (G3). Squaring \emph{before}
integrating avoids the FMP bottleneck.
\end{remark}

\section{Conditionality}\label{sec:conditionality}

\textbf{Inputs used after \textup{($\Pi$-int)} and the auxiliary annular/perimeter cap:}
slicing identity (Maggi); bad-ray multiplicity bound (elementary, see proof of Lemma~\ref{lem:T1rad}); slicewise good-ray equality; spherical Fubini;
divergence identity (5.3); Talenti profile gap;
slice-rearrangement Lemma \ref{lem:slice-rearrange}; (F) of
\cite{CentroidBound}.

\textbf{Inputs explicitly not used:} the per-$\rho$ pointwise radial-trace inputs (1.1), (1.2)=(R),
(2.4); pointwise (B²); the boundary-layer pointwise (7.1) of
\cite{BoundaryLayerTransfer}; any selection, graph entry,
Schauder, Fuglede.

\textbf{Centroid-route status:} Theorem~\ref{thm:PiInput}, plus the
annular/perimeter cap stated in Theorem~\ref{thm:T1sq}, would complete this
centroid subroute. Without them the present block proves the exact
space--time identity and a conditional centroid closure only. The repaired
metric closure is instead supplied by \cite{WeightedMetricTrace}.

\begin{remark}[How the metric closure uses $D_H$ and $D_I$ jointly]
\label{rem:DH-DI-joint}
The repaired metric route \cite{WeightedMetricTrace} bypasses
Theorem~\ref{thm:PiInput} (and hence the centroid space--time control
of physical crossings) by pairing the two pieces of the deficit
identity $D_H+D_I$ of \cite{LevelSetIdentity} as two independent
integral inputs:

\begin{itemize}
\item the integrated isoperimetric defect
$\int(P(E_\rho)-P(B_\rho))(P(E_\rho)+P(B_\rho))\,d\mu\le C\delta_T$
(i.e.\ $\int D_I\,d\mu\le C\delta_T$) controls the $L^2(d\mu)$-norm of
the levelwise Fraenkel distance via Fusco--Julin
\cite[Lemma~3.1]{WeightedMetricTrace};

\item the integrated Cauchy--Schwarz defect
$\int D_H\,d\mu\le C\delta_T$ controls the $L^2(d\mu)$-norm of the
metric derivative $|\dot F_\rho|_{\mathcal X}$ in the quotient space
$\mathcal X$ of \cite{MetricFramework} via its metric first variation.
\end{itemize}

The first input says ``many levels are close to balls'' but does not
fix \emph{which} level; the second says ``the family $\rho\mapsto E_\rho$
does not curve away from its trajectory faster than $\sqrt{\delta_T}$
in $L^2$''.  Combining them on the order-one window
$J^*=[\rho_\delta-L^*,\rho_\delta]$ via Markov (on the distance term) +
Cauchy--Schwarz kinetic transport (on the velocity term) yields a
\emph{uniform} pointwise bound $d_{\mathcal X}(F_\rho,B_1)^2\le C_*\delta_T$ on
$J^*$, including the endpoint $\rho=\rho_\delta$ where the
boundary-layer transfer of \cite{BoundaryLayerTransfer} has minimal
thickness $1-\rho_\delta\asymp\sqrt{\delta_T}$.  This is the
$D_H/D_I$-route to sharp Faber--Krahn stability: each integrated
defect supplies a different one of the two $H^1$-in-$\rho$ ingredients
required by the 1D Sobolev embedding underlying the pointwise endpoint
bound, while neither alone would suffice. The centroid space--time
$\Pi$-control of Theorem~\ref{thm:PiInput} is logically independent of
this argument and would only be needed for the alternative
centroid-route closure.
\end{remark}

\begin{thebibliography}{99}

\bibitem{LevelSetIdentity}
Plan~2 building block \textsc{LevelSetIdentity}.

\bibitem{MetricFramework}
Plan~2 building block \textsc{MetricFramework}.

\bibitem{CentroidBound}
Plan~2 building block \textsc{CentroidBound}.

\bibitem{WeightedMetricTrace}
Plan~2 building block \textsc{WeightedMetricTrace}.

\bibitem{BoundaryLayerTransfer}
Plan~2 building block \textsc{BoundaryLayerTransfer}.

\bibitem{Maggi2012}
F.~Maggi, Cambridge Univ. Press, 2012. Thm 15.9, 18.11, 19.4.

\end{thebibliography}

\end{document}
